\scanpage{62}
\uline{La fonction $\unsure{p|\wp}z$}

Définition intrinsèque : c'est une fonction elliptique d'ordre 2, admettant l'origine pour pole double, \deleted{\unsure{y ayant pour pa}} \inserted{la fonction} $pz-\dfrac{1}{z^2}$ s'annulant à l'origine. Si $\omega_1$ et $\omega_2$ sont donnés, cela détermine sans plus cette fonction elliptique.

Elle est paire [car $\varphi(z)-\varphi(-z)$ est elliptique, n'a pas de pole, s'annule à l'origine] \deleted{et vérifie l'équation} \inserted{et est donc de la forme}
\[
\deleted{\unsure{p^{\prime2}z}}\qquad
pz=\frac1{z^2}+c_1z^{\unsure{2|3}}+c_{\unsure{1|2}}z^4\mathord{\text{---}}
\]

Elle vérifie l'équation
\[
p^{\prime 2}z=4p^3z-g_2pz-g_3\qquad
g_1=20c_1\qquad g_3=28c_2
\]

C'est donc la fonction inverse de
\[
z=\int_\infty^Z\frac{dZ}{\sqrt{4Z^3-g_2Z-g_3}}
\]

On en déduit que si on pose
\[
p\omega_1=e_1\qquad p\omega_2=e_2\qquad p\omega_3=e_3\ ,
\]
$e_1,e_2,e_3$ sont les racines de l'équation $4Z^3-g_2Z-g_3=0$ ; on a donc $p'\omega_1=p'\omega_2=p'\omega_3=0$

On a \deleted{$p'\omega_1=\unsure{e_1}$}
\[
\begin{aligned}
p''z&=6p^{\unsure{2|\prime}}z-\frac{g_2}{2}\\
p^{\unsure{\prime\prime\prime|\prime\prime}}z&=12pz\,p'z\\
&\mathord{-}\quad\mathord{-}\quad\mathord{-}\quad\mathord{-}
\end{aligned}
\]

$p^{\unsure{(2n)|(2n-1)}}z$ s'exprime en fonction de $pz$ par un polynome de degré $n+1$ ; $p^{2n+1)}z$ par le produit par $p'z$ d'un polynome en $pz$ de degré $n$.

\inserted{Développ\textsuperscript{mt} en série double}

\uline{La fonction $\zeta z$}. S'introduit par l'intégration des fonctions elliptiques.

La fonction $\int_0^z(pz-\frac1{z^2})dz$ est uniforme, en lui ajoutant $\frac1z$ on obtient une primitive de $\unsure{pz|\wp z}$. On pose
\[
\zeta z=\frac1z-\int_0^z(pz-\frac1{z^2})dz\qquad\text{On a }\quad\zeta'z=-pz
\]
% archive sheet number: 61
